\scanpage{7}
\unsure{générisation|généralisation}. Pour achever de prouver la prop, il suffit de prouver (compte tenu de \textcircled{\scriptsize i} \inserted{et \textcircled{\scriptsize ii}}) que pour tout $\alpha\in F(M)$, $M$ de t.f. sur $A$, il existe un ``module de définition $N$'' i.e un plus petit sous-module $N$ de $M$ tel que $\alpha\in\operatorname{Im}(F(N)\to F(M))$. Or on a pour tt $p\in S$ un $N_p\subset M_p$ $(=\operatorname{Im}(\alpha_p:P_p\unsure{-|\to} M_p))$. \deleted{Soit tt} Il suffit de voir que le système des $N_p$ provient d'un sous-module $N$ de $M$ (on voit alors aisément qu'il conviendra). Soit $U\subset S$ le plus grand ouvert sur lequel ça marche, \inserted{par un $\underline{\widetilde N}_U\subset\widetilde M|U$} Si on avait $U\ne S$, soit $x\in S-U$ \deleted{\unsure{un pt maxi|complètement}} un pt maximal de $S-U$.

\deleted{Il existe un voisinage \unsure{ouvert}} \deleted{Il existe un} Soit $S'=\operatorname{Spec}\mathcal O_{S,x}$, $U'=U\cap S'$, \deleted{donc} $N_{S'}=\widetilde{N_x}$ (NB. \tikz[baseline=(a.base)]{\node[inner sep=1pt](a){$x=p$};\draw (a.center) ellipse (.48cm and .23cm);} !), alors $\widetilde N_x|U'=\underline N_U|U'$, donc $\exists$ voisinage ouvert $V$ de $x$, \deleted{tel que} et un $\underline{N}'\subset\widetilde M|V$, tel que $\underline N'_x=N_x\subset M_x$ et $\underline N'|V\cap U=\underline N_U|V$, i.e.
\[
\left\{\begin{aligned}
\underline N'_x&=N_x\subset M_x\\
\underline N'_s&=N_s\subset M_s\qquad\text{si }\deleted{\unsure{x}}s\in U\cap V.
\end{aligned}\right.
\]
\margin{$\alpha_V\in F(N')$ donc $N_s\subset\underline N'_s$ pour tout $s\in V$}
\margin{$\longrightarrow$ Donc on est ramené au cas $\underline N'=\unsure{\widetilde M|M}$ \deleted{\unsure{quitte à}}}
Prenant $V$\inserted{$=\operatorname{Spec} A_f$} assez petit, \deleted{la condition (v) implique} \inserted{on peut supposer que} \deleted{\unsure{et}} la condition (v) implique qu'on peut supposer que $F_{(A/p)_f}$ est représentable\unsure{?|.}\ \deleted{Il s'ensuit que} \deleted{et} On voit tout de suite que ceci implique que (quitte à restreindre encore $V$) l'image de $N_s$ dans $M_s$ \deleted{est \uline{égale} \unsure{déjà}} \inserted{est tout} modulo $p$, donc est tout pour $s\in V(p)$ par Nakayama. \unsure{Ceci gagne|Ceci signifie|C'est gagné|Cela suffit}.

% archive sheet number: 6
